\section{Exercise 1.} Show that $\mathcal{E}_P(g) = R_P(g) - R_P(g^*) = \EE[|g(X)-g^*(X)|\cdot|2\eta^*(X)-1|]$, for the binary classification setting.

\textit{Solution.} For every $g\in\mathcal{G}$, we have
\begin{align*}
  R_P(g) & = \EE[\mathds{1}(Y\ne g(X))] = \EE[\EE[\mathds{1}(Y\ne g(X)) | X]].
\end{align*}
Define the pointwise risk $r_g:\mathcal{X}\to [0,1]$ by
\begin{align*}
  r_g(x) & =\EE[\mathds{1}(Y\ne g(X)) | X=x]  = \PP(Y\ne g(X) | X=x) = \PP(Y= 1-g(x)| X=x) \\
         & = \begin{cases}
               \eta^*(x), \text{ if } g(x)=0 \\
               1- \eta^*(x), \text{ if } g(x)=1.
             \end{cases}
\end{align*}
Let us compute $r_g(x) - r_g^*(x)$ for each $x\in\mathcal{X}$. Consider the following cases.
\begin{enumerate}
  \item If $g(x)=g^*(x)$, then $r_g(x) - r_g^*(x)=0$.
  \item If $g(x)\ne g^*(x)$
        \begin{enumerate}
          \item [(a)] If $g^*(x) = 0$, then $\eta^*(x)<\dfrac{1}{2}$ and $r_{g^*}(x)=\eta^*(x)$. Also, $r_g(x)=1-\eta^*(x)$, since $g(x) = 1$. Hence, $r_g(x) - r_g^*(x)=1-2\eta^*(x)$.
          \item [(b)] If $g^*(x) = 1$, then $\eta^*(x)\ge\dfrac{1}{2}$ and $r_{g^*}(x)=1-\eta^*(x)$. Also, $r_g(x)=\eta^*(x)$, since $g(x) = 0$. Hence, $r_g(x) - r_g^*(x)=2\eta^*(x)-1$.
        \end{enumerate}
\end{enumerate}
Notice that in all cases, $r_g(x) - r_g^*(x)\ge 0$. This verifies that $g^*$ is indeed optimal. Collecting all cases, we write
$$r_g(x) - r_g^*(x) = |g(x)-g^*(x)|\cdot|2\eta^*(x)-1|.$$ Therefore,
$$\E_P(g) = \EE[r_g(X)-r_{g^*}(X)] = \EE[|g(X)-g^*(X)|\cdot|2\eta^*(X)-1|].$$

\section{Exercise 2.} Show that $\E_P(g) = \EE[(g(X)-g^*(X))^2]$, for the regression setting.

\textit{Solution.} Similarly to Exercise 1, we use the pointwise risk $r_g:\X\to \RR_+$. We have
\begin{align*}
  r_g(x) & = \EE[(Y-g(X))^2 | X=x]                                                                    \\
         & = \EE[(Y-g^*(X)+g^*(X)-g(X))^2 | X=x]                                                      \\
         & = \EE[(Y-g^*(X))^2 + (g^*(X)-g(X))^2 + 2\langle Y-g^*(X),g^*(X)-g(X)\rangle | X=x]         \\
         & = \EE[(Y-g^*(X))^2|X=x] + (g^*(x)-g(x))^2 + 2\langle \EE[Y-g^*(X)|X=x],g^*(x)-g(x)\rangle.
\end{align*}
We have $\EE[(Y-g^*(X))^2|X=x]=r_{g^*}(x)$. Also, $\EE[Y-g^*(X)|X=x]= \EE[Y|X=x] - g^*(x) = 0$ by the definition of $g^*(x)$. Hence, $r_g(x) = r_{g^*}(x) + (g^*(x)-g(x))^2$, or
$$r_g(x) - r_{g^*}(x) = (g^*(x)-g(x))^2,$$
which yields the desired result.

\section{Exercise 3} Provide an upper bound of $R_P(\hat{g}_{n,\G}) - R_P(g^*_{\G})$ by the VC dimension.

\textit{Solution.}

\textbf{Step 1.} The first step is to derive the Vapnik's Lemma to get rid of $\hat{g}_{n,\G}$ and $g^*_{\G}$. We have
\begin{align*}
  R_P(\hat{g}_{n,\G}) - R_P(g^*_{\G})
   & = R_P(\hat{g}_{n,\G}) - \hat{R}_n(\hat{g}_{n,\G}) + \underbrace{\hat{R}_n(\hat{g}_{n,\G}) -  \hat{R}_n(g^*_{\G})}_{\le 0} + \hat{R}_n(g^*_{\G}) - R_P(g^*_{\G}) \\
   & \le 2\sup_{g\in\G}|R_P(g) - \hat{R}_n(g)|.
\end{align*}

\textbf{Step 2.} Next, we use McDiarmid's inequality to pass the random $\sup_{g\in\G}|R_P(g) - \hat{R}_n(g)|$ to its expectation. Define $f:\X\times Y \to \RR_+$ by
$$f(z_1,\ldots,z_n) = \sup_{g\in\G}|R_P(g) - \dfrac{1}{n}\sum_{i=1}^n \mathds{1}(y_i\ne g(x_i))|,$$
where $z_i=(x_i,y_i)$ for each $i\in\{1,\ldots,n\}$. Then, for each $i\in\{1,\ldots,n\}$, we have
\begin{align*}
  |f(z_1,\ldots,z_i,\ldots,z_n) - f(z_1,\ldots,z_i',\ldots,z_n)| & \le \sup_{g\in\G}\left|\dfrac{1}{n}\sum_{i=1}^n \mathds{1}(y_i\ne x_i) - \dfrac{1}{n}\sum_{i=1}^n \mathds{1}(y_i'\ne x_i)\right| \\
                                                                 & = \dfrac{1}{n}\sup_{g\in\G}|\mathds{1}(y_i\ne x_i) - \mathds{1}(y_i'\ne x_i)|                                                    \\
                                                                 & \le \dfrac{1}{n}.
\end{align*}
Hence $f$ has bounded differences with $c_i = \dfrac{1}{n}$ for each $i\in\{1,\ldots,n\}$. Let $Z=f(Z_1,\ldots, Z_n)$, then by McDiarmid's inequality, for every $t\in(0,1)$, with probability at least $1-t$,
$$\PP(Z- \EE[Z] \ge t) \le \exp\left(-\frac{2t^2}{\sum_{i=1}^n c_i^2}\right) = \exp(-2nt^2).$$
Letting $\exp(-2nt^2) = \delta$, we have $t = \sqrt{\dfrac{\log(1/\delta)}{2n}}$. Hence, with probability at least $1-\delta$,
$$\sup_{g\in\G}|R_P(g) - \hat{R}_n(g)| \le \EE\left[\sup_{g\in\G}|R_P(g) - \hat{R}_n(g)|\right] + \sqrt{\dfrac{\log(1/\delta)}{2n}}.$$

\textbf{Step 3.} We use symmetrization to get rid of the unknown $R_P$. Introduce independent samples $Z'_1,\ldots,Z'_n$ with the same distribution as $Z=(Z_1,\ldots,Z_n)$. Let $\hat{R}'_n(g)=\dfrac{1}{n}\mathds{1}_{Y'_i\ne g(X'_i)}$. Then,
$$\EE\left[\left.\hat{R}'_n(g)\right| Z\right] = \EE\left[\hat{R}'_n(g) \right]=\EE[\hat{R}_n(g)] = R_P(g).$$
By Jensen's inequality,
\begin{align*}
  |R_P(g) - \hat{R}_n(g)|
   & = \left|\EE\left[\left.\hat{R}'_n(g)-\hat{R}_n(g)\right| Z\right]\right| \le \EE\left[\left.\left|\hat{R}'_n(g)-\hat{R}_n(g)\right|\,\,\right|Z\right].
\end{align*}
Taking the supremum of both sides over $g\in\G$, we get
$$\sup_{g\in\G}|R_P(g) - \hat{R}_n(g)| \le \sup_{g\in\G}\EE\left[\left.|\hat{R}'_n(g)-\hat{R}_n(g)|\,\,\right|Z\right] \le \EE\left[\sup_{g\in\G}\left.|\hat{R}'_n(g)-\hat{R}_n(g)|\,\,\right|Z\right].$$
Taking the expectation of both sides, we get
\begin{align*}
  \EE\left[\sup_{g\in\G}|R_P(g) - \hat{R}_n(g)|\right]
   & \le \EE\left[\sup_{g\in\G}|\hat{R}'_n(g)-\hat{R}_n(g)|\right]                                                                    \\
   & = \EE\left[\sup_{g\in\G}\left|\dfrac{1}{n}\sum_{i=1}^n (\mathds{1}(Y'_i \ne g(X'_i)) - \mathds{1}(Y_i \ne g(X_i)))\right|\right]
\end{align*}

Introduce Rademacher random variables $\sigma_1,\ldots,\sigma_n$ independent of $Z$ and $Z'$. Then $\mathds{1}(Y'_i \ne g(X'_i)) - \mathds{1}(Y_i \ne g(X_i))$ and $\sigma_i(\mathds{1}(Y'_i \ne g(X'_i)) - \mathds{1}(Y_i \ne g(X_i)))$ have the same distribution. Hence,
\begin{align*}
  \EE\left[\sup_{g\in\G}|R_P(g) - \hat{R}_n(g)|\right]
   & \le \EE\left[\sup_{g\in\G}\left|\dfrac{1}{n}\sum_{i=1}^n \sigma_i(\mathds{1}(Y'_i \ne g(X'_i)) - \mathds{1}(Y_i \ne g(X_i)))\right|\right] \\
   & \le 2\EE\left[\sup_{g\in\G}\left|\dfrac{1}{n}\sum_{i=1}^n \sigma_i\mathds{1}(Y_i \ne g(X_i))\right|\right]
\end{align*}

\textbf{Step 4.} Push the randomness to the Rademacher variables only.

\begin{align*}
  \EE\left[\sup_{g\in\G}\left|\dfrac{1}{n}\sum_{i=1}^n \sigma_i\mathds{1}(Y_i \ne g(X_i))\right|\right]       \le \sup_{(x,y)\in (\X\times Y)^n} \EE\left[ \sup_{g\in\G}\left|\dfrac{1}{n}\sum_{i=1}^n \sigma_i\mathds{1}(y_i \ne g(x_i))\right|\right] \\
  = \sup_{(x,y)\in (\X\times Y)^n} \EE\left[\dfrac{1}{n} \sup_{g\in\G}\sum_{i=1}^n \sigma_i\mathds{1}(y_i \ne g(x_i))\right]
\end{align*}

Note that we can remove the absolute value since the Rademacher variables are symmetric.

\textbf{Step 5.} Relate the last term to the VC dimension, then use Massart's Lemma, since the lemma allows us to work directly with expectation, without passing back to probability like Hoeffding's inequality. Since $g\mapsto \sum\limits_{i=1}^{n}\mathds{1}(y_i\ne g(x_i))$ can only take at most $V_{\G}(n)$ values. Hence, there exists $g_1,\ldots,g_{V_{\G}(n)}\in\G$ such that
\begin{align*}
  \EE\left[\dfrac{1}{n} \sup_{g\in\G}\sum_{i=1}^n \sigma_i\mathds{1}(y_i \ne g(x_i))\right] = \EE\left[\dfrac{1}{n} \max_{j\in\{1,\ldots,V_{\G}(n)\}}\sum_{i=1}^n \sigma_i\mathds{1}(y_i \ne g_j(x_i))\right] \le \sqrt{\dfrac{2\log V_\G(n)}{n}}.
\end{align*}
Taking the supremum over $(x,y)\in (\X\times Y)^n$ yields the same bound, since the bound does not depend on $(x,y)$.

\textbf{Step 6.} By Sauer's lemma, $V_\G(n) \le \left(\dfrac{en}{d}\right)^d$ where $d$ is the VC dimension of $\G$. Collecting all steps, we have that with probability at least $1-\delta$,
\begin{align*}
  R_P(\hat{g}_{n,\G}) - R_P(g^*_{\G}) \le 4 \sqrt{\dfrac{2d\log(en/d)}{n}} + \sqrt{\dfrac{2\log(1/\delta)}{n}}.
\end{align*}

